\en{We prove (using App.~\ref{\en{EN}kapdivi}) that $\sqrt{D}\cdot\frac{E_2^*(\tau)}{\eta^4(\tau)}\cdot(AC)^2$ is an algebraic integer if $\tau$ satisfies $C\tau^2+B\tau+A=0$ with discriminant $D$.}%
\de{Wir beweisen, dass $\sqrt{D}\cdot\frac{E_2^*(\tau)}{\eta^4(\tau)}\cdot(AC)^2$ ganzalgebraisch ist, falls $\tau$ eine Lösung von $C\tau^2+B\tau+A=0$ mit Diskriminante $D$ ist.}
\en{The proof elaborates \cite[Lemma A3]{\en{EN}Masser1975}.}%
\de{Der Beweis arbeitet \cite[Lem.~A3]{\en{EN}Masser1975} aus.}

\begin{defi}\label{\en{EN}bemkonv}
\en{Throughout this appendix, the following notations are used:}%
\de{In diesem ganzen Anhang gelten folgende Konventionen:}
\begin{itemize}
    \item \en{We fix a lattice $L=\mathbb Z \omega_1 + \mathbb Z \omega_2$ with complex multiplication as in Def.~\ref{\en{EN}defiCM},
    where $\tau$ denotes the period ratio $\tau:=\frac{\omega_2}{\omega_1}\in \mathbb H$ and \ldots}%
    \de{Wir betrachten ein festes Gitter $L=\mathbb Z \omega_1 + \mathbb Z \omega_2$ mit komplexer Multiplikation wie in Def.~\ref{\en{EN}defiCM},
    wobei wir das Periodenverhältnis mit $\tau:=\frac{\omega_2}{\omega_1}\in \mathbb H$ bezeichnen und \ldots}
    \begin{itemize}
    \item \en{$A$, $B$ and $C$ are integers with $\operatorname{gcd}(A;B;C)=1$ and $A+B\tau+C\tau^2=0$,}%
    \de{$A$, $B$ und $C$ seien ganze Zahlen mit $\operatorname{ggT}(A;B;C)=1$ und $A+B\tau+C\tau^2=0$,}
    \item \en{and $D$ denotes the discriminant of $\tau$, which is $D = B^2-4AC$.}%
    \de{und $D$ bezeichne die Diskriminante von $\tau$, also $D = B^2-4AC$.}
    \end{itemize}
    \item \en{As in Def.~\ref{\en{EN}defis2}, let $E_2^*(\tau):=E_2(\tau) - \frac{3}{\pi Im(\tau)}$ with $E_2$ from Thm.~\ref{\en{EN}fouriertheorem}.}%
    \de{Wie in Def.~\ref{\en{EN}defis2} sei $E_2^*(\tau):=E_2(\tau) - \frac{3}{\pi Im(\tau)}$ mit $E_2$ aus Thm.~\ref{\en{EN}fouriertheorem}.}
    \item \en{$\eta_1$ and $\eta_2$ denote the basic quasi periods of $L$ (cf.~Def.~\ref{\en{EN}defetak}).}%
    \de{$\eta_1$ und $\eta_2$ bezeichnen die Basis-Quasiperioden von $L$ (vgl.~Def.~\ref{\en{EN}defetak}).}
    \item \en{$\zeta(z)$ and $\wp(z)$ denote the Weierstraß functions from Def.~\ref{\en{EN}defizeta} and~\ref{\en{EN}defiwp}.}%
    \de{$\zeta(z)$ und $\wp(z)$ bezeichnen die Weierstraß'schen Funktionen aus Def.~\ref{\en{EN}defizeta} und~\ref{\en{EN}defiwp}.}
    \item \en{$\DIV(m)$ denotes the set of $m$-division points in the fundamental parallelogram $\mathcal P =\left\{~s\omega_1+t\omega_2~|~0\leq s,t < 1~\right\}$ as in Def.~\ref{\en{EN}defidivi}.}%
    \de{$\DIV(m)$ bezeichne die Menge der $m$-Teilungsstellen im Periodenparallelogramm $\mathcal P =\left\{~s\omega_1+t\omega_2~|~0\leq s,t < 1~\right\}$ wie in Def.~\ref{\en{EN}defidivi}.}
\end{itemize}
\end{defi}

\begin{defi}\label{\en{EN}defkappa}
\en{In the notations of Def.~\ref{\en{EN}bemkonv}, we define $\kappa$ by}%
\de{In der Notation von Def.~\ref{\en{EN}bemkonv} definieren wir $\kappa$ durch}
$$\kappa\omega_2 := A\eta_1 - C\tau\eta_2$$
\end{defi}

\begin{lem}\label{\en{EN}lemkapa1}
\en{In the notations of Def.~\ref{\en{EN}bemkonv}, it holds: }%
\de{In der Notation von Def.~\ref{\en{EN}bemkonv} gilt:}
$$\kappa = -\sqrt{D}\cdot\frac{\pi^2}{3\omega_1^2}\cdot E_2^*(\tau)$$
\end{lem}
\begin{proof}
\en{First we multiply Def.~\ref{\en{EN}defkappa} with $\omega_1$ and obtain:}%
\de{Zunächst multiplizieren wir Def.~\ref{\en{EN}defkappa} mit $\omega_1$ und erhalten:}
$$\kappa\omega_1\omega_2 = A\omega_1\eta_1 - C \tau \omega_1\eta_2$$
\en{Then we use Legendre's relation (Prop.~\ref{\en{EN}legendre}) and obtain:}%
\de{Dann nutzen wir die Legendre'sche Relation (Satz~\ref{\en{EN}legendre}):}%
\begin{align*}
    \kappa\omega_1\omega_2 = A\omega_1\eta_1 - C \tau (\eta_1\omega_2-2\pi i)
    &=\left(A-C\tau\frac{\omega_2}{\omega_1}\right)\cdot \omega_1\eta_1 + 2\pi i C \tau\\
    &=\left(A-C\tau^2\right)\cdot \omega_1\eta_1 + 2\pi i C \tau
\end{align*}
\en{Thm.~\ref{\en{EN}fouriertheorem} tells $\eta_1(L_\tau)=\frac{\pi^2}{3}\cdot E_2(\tau)$.
Now we have $L=\omega_1\cdot L_\tau$ and from Prop.~\ref{\en{EN}etatransf} we get $\eta_1(L)=\frac{1}{\omega_1}\cdot\eta_1(L_\tau)=\frac{1}{\omega_1}\cdot \frac{\pi^2}{3}\cdot E_2(\tau)$.
Thus $\omega_1\eta_1 = \frac{\pi^2}{3}\cdot E_2(\tau)$ and we obtain:}%
\de{In Thm.~\ref{\en{EN}fouriertheorem} haben wir $\eta_1(L_\tau)=\frac{\pi^2}{3}\cdot E_2(\tau)$ bewiesen.
Jetzt gilt $L=\omega_1\cdot L_\tau$ und Satz~\ref{\en{EN}etatransf} liefert $\eta_1(L)=\frac{1}{\omega_1}\cdot\eta_1(L_\tau)=\frac{1}{\omega_1}\cdot \frac{\pi^2}{3}\cdot E_2(\tau)$.
Somit folgt $\omega_1\eta_1 = \frac{\pi^2}{3}\cdot E_2(\tau)$ und}
\begin{align*}
    \kappa\omega_1\omega_2 &=\left(A-C\tau^2\right)\cdot \frac{\pi^2}{3}\cdot E_2(\tau) + 2\pi i C \tau
\end{align*}
\en{From}\de{Aus} $A+B\tau=-C\tau^2$ \en{we obtain}\de{folgt} $A-C\tau^2 = A+(A+B\tau)=2A+B\tau$ \en{and}\de{und}
\begin{align*}
   - \sqrt{D}\cdot \tau &= -\sqrt{D}\cdot\frac{-B+\sqrt{D}}{2C} = \frac{B\cdot\sqrt{D}-D}{2C} = \frac{B\cdot\sqrt{D}-B^2+4AC}{2C}\\
   &= 2A + \frac{B\cdot\sqrt{D}-B^2}{2C} = 2A +B\tau = A-C\tau^2
\end{align*}
\en{which yields}\de{Das führt zu}{\abovedisplayskip=-4pt
\begin{align*}
    \kappa\omega_1\omega_2 &=- \sqrt{D}\cdot \tau\cdot \frac{\pi^2}{3}\cdot E_2(\tau) + 2\pi i C \tau\\
    &= - \sqrt{D}\cdot \tau\cdot \frac{\pi^2}{3}\cdot \left(E_2(\tau) - \frac{2\pi i C \tau}{\sqrt{D}\cdot \tau\cdot \frac{\pi^2}{3}}\right)
\end{align*}}
\en{But from}\de{Aber aus} $\im(\tau)=\frac{\sqrt{-D}}{2C}= \frac{\sqrt{D}}{2C\cdot i}$ \en{we get}\de{folgt nun}
\begin{align*}
    \kappa\omega_1\omega_2 &= - \sqrt{D}\cdot \tau\cdot \frac{\pi^2}{3}\cdot \left(E_2(\tau) - \frac{3}{\pi\im(\tau)}\right)
\end{align*}
\en{Finally we divide by}\de{Schließlich dividieren wir durch} $\omega_1\cdot\omega_2$ \en{and obtain}\de{und erhalten} $\kappa = - \sqrt{D}\cdot \frac{\tau}{\omega_1\omega_2}\cdot \frac{\pi^2}{3}\cdot E_2^*(\tau)$.
\en{Here we simplify}\de{Hier vereinfachen wir noch} $\frac{\tau}{\omega_1\omega_2}=\frac{\omega_2/\omega_1}{\omega_1\omega_2}=\frac{1}{\omega_1^2}$ \en{which proves the Lemma.}\de{womit das Lemma bewiesen ist.}
\end{proof}

\pagebreak
\begin{lem}\label{\en{EN}lemfellip}
\en{In the notations of Def.~\ref{\en{EN}bemkonv}, the function}%
\de{In der Notation von Def.~\ref{\en{EN}bemkonv} ist die Funktion}
$$f(z):=-A\zeta(Cz)+C\tau\zeta(C\tau z) + C\tau\kappa z$$
\en{is elliptic with periods $\omega_1$ and $\omega_2$.}%
\de{elliptisch mit Perioden $\omega_1$ und $\omega_2$.}
\end{lem}
\begin{proof}
\en{To prove the first period we calculate}\de{Um die erste Periode zu beweisen, berechnen wir}
\begin{align*}
    f(z+\omega_1)-f(z) &= -A\cdot\underbrace{\left(\zeta(C(z+\omega_1))-\zeta(Cz)\right)}_{T_1}\\
    &~~~+ C \tau \cdot \underbrace{\left(\zeta(C\tau(z+\omega_1))-\zeta(C\tau z)\right)}_{T_2} +~ C\tau\kappa\omega_1
\end{align*}
\en{For calculating $T_1 = \zeta(Cz+C\omega_1)-\zeta(Cz)$ we apply $C$ times the Definition~\ref{\en{EN}defetak} of $\eta_1=\zeta(Cz+\omega_1)-\zeta(Cz)$ and get
$T_1=C\eta_1$.
Similarly, as $C\tau\omega_1=C\omega_2$, we get $T_2 = \zeta(C\tau z+C\tau\omega_1)-\zeta(C\tau z) = C\eta_2$. This yields}%
\de{Um den Wert von $T_1 = \zeta(Cz+C\omega_1)-\zeta(Cz)$ zu bestimmen, wenden wir $C$ mal die Def.~\ref{\en{EN}defetak} von $\eta_1=\zeta(Cz+\omega_1)-\zeta(Cz)$ an und erhalten
$T_1=C\eta_1$.
Aus $C\tau\omega_1=C\omega_2$ folgt genauso $T_2 = \zeta(C\tau z+C\tau\omega_1)-\zeta(C\tau z) = C\eta_2$. Dies führt zu}
\begin{align*}
    f(z+\omega_1)-f(z) &= - A\cdot C \eta_1 + C\tau \cdot C\eta_2 + C\kappa\omega_2 = C\cdot\left(\kappa\omega_2-A\eta_1+C\tau\eta_2\right)=0,
\end{align*}
\en{where we used Def.~\ref{\en{EN}defkappa} of $\kappa$ in the last step.}%
\de{wobei wir im letzten Schritt Def.~\ref{\en{EN}defkappa} von $\kappa$ benutzt haben.}
\en{Now the second period:}\de{Nun zur zweiten Periode:}
\begin{align*}
    f(z+\omega_2)-f(z) &= -A\cdot\underbrace{\left(\zeta(C(z+\omega_2))-\zeta(Cz)\right)}_{T_3}\\
    &~~~+ C \tau \cdot \underbrace{\left(\zeta(C\tau(z+\omega_2))-\zeta(C\tau z)\right)}_{T_4}
    +~C\tau\kappa\omega_2
\end{align*}
\en{As above we have $T_3=C\eta_2$. Then it holds}%
\de{Wie oben gilt $T_3=C\eta_2$. Außerdem ist}
$$C\tau\omega_2=C\tau^2\omega_1 = -(A+B\tau)\omega_1 = -A\omega_1-B\omega_2$$
\en{which yields $T_4 = -A\eta_1-B\eta_2$. This shows}%
\de{was $T_4=-A\eta_1-B\eta_2$ beweist und somit}
\begin{align*}
    f(z+\omega_2)-f(z) &= - A\cdot C \eta_2 + C\tau \cdot (-A\eta_1-B\eta_2) + C\tau\kappa\omega_2\\
    &= -AC\eta_2-AC\tau\eta_1-BC\tau\eta_2+C\tau(A\eta_1-C\tau\eta_2)\\
    &= -C\cdot\eta_2\cdot\left(A+B\tau+C\tau^2\right)=0
\end{align*}
\en{Thus the function $f(z)$ has the periods $\omega_1$ and $\omega_2$.}%
\de{Also hat die Funktion $f(z)$ die Perioden $\omega_1$ und $\omega_2$.}
\end{proof}

\begin{lem}\label{\en{EN}lemctau}
\en{In the notations of Def.~\ref{\en{EN}bemkonv}, it holds:}%
\de{In der Notation von Def.~\ref{\en{EN}bemkonv} gilt:}
\en{The number of $C\tau$-division points in $\mathcal P$ is $AC-1$. And every $C\tau$-division-point is also an $AC$-division-point. In brief:}%
\de{Die Anzahl der $C\tau$-Teilungsstellen in $\mathcal P$ ist $AC-1$. Und jede $C\tau$-Teilungsstelle ist auch eine $AC$-Teilungsstelle. Kurz:}
$$|\DIV(C\tau)|=AC-1\qquad\text{\en{and}\de{und}}\qquad \DIV(C\tau)\subset \DIV(AC)$$
\end{lem}
\begin{proof}\en{Let}\de{Sei} $L=\mathbb Z \omega_1 + \mathbb Z \omega_2$
\en{with}\de{mit} $\tau:=\frac{\omega_2}{\omega_1}\in CM$
\en{as in}\de{wie in} Def.~\ref{\en{EN}bemkonv}.\\
\en{Then it holds for all $u\in \DIV(C\tau)$:}\de{Dann gilt für alle $u\in \DIV(C\tau)$:}
\begin{align*}
    C\tau u \in L
    \quad\Longleftrightarrow\quad
     C\tau u &= -k\omega_1 + l \omega_2
    \quad\Longleftrightarrow\quad
    u = \frac{-k\omega_1 + l \omega_2}{C\tau} = \frac{-k\omega_1\bar\tau + l \omega_2\bar\tau}{C\tau\bar\tau}
\end{align*}
\en{Then we have $\tau\bar\tau = \frac{B^2-(B^2-4AC)}{4C^2} = \frac A C$ and $\tau+\bar\tau = -\frac B C$.\\
This leads to $\bar\tau = -\frac B C-\tau$ and $\omega_2\bar\tau=\omega_1\tau\bar\tau=\omega_1\cdot\frac{A}{C}$ and $\omega_1\tau=\omega_2$, thus:}%
\de{Weiter gilt $\tau\bar\tau = \frac{B^2-(B^2-4AC)}{4C^2} = \frac A C$ und $\tau+\bar\tau = -\frac B C$.\\
Hieraus folgt $\bar\tau = -\frac B C-\tau$ und $\omega_2\bar\tau=\omega_1\tau\bar\tau=\omega_1\cdot\frac{A}{C}$ sowie $\omega_1\tau=\omega_2$, also:}
\begin{align}
    u &= \frac{-k\cdot\omega_1\cdot\left(-\frac{B}{C}-\tau\right)+l\cdot\omega_1\cdot A/C}{C\cdot A/C} = \frac{k}{A}\cdot\omega_2+\frac{l\cdot A+k\cdot B}{AC}\cdot\omega_1\label{\en{EN}ctauu}
\end{align}
\en{Next, $u$ has to be in $\mathcal P$ which yields $0\leq \frac k A < 1$ and $0\leq \frac{l\cdot A+k\cdot B}{AC} <1$.
From the first we see that there are $A$ possible values of $k$. From the second we see that $0\leq l + k\cdot\frac{B}{A} < C$, thus we have $C$ possible values of $l$, no matter what value of $k$ we had chosen. This yields $AC$ values of $u$ in $\mathcal P$ such that $C\tau u\in L$. Since $u=0$ is not a $C\tau$-division-point, we have $AC-1$ of the $C\tau$-division-points in $\mathcal P$.}%
\de{Außerdem soll $u$ in $\mathcal P$ liegen, also muss $0\leq \frac k A < 1$ und $0\leq \frac{l\cdot A+k\cdot B}{AC} <1$ gelten.
Aus der ersten Ungleichung folgt, dass es $A$ mögliche Werte von $k$ gibt. Aus der zweiten erkennen wir $0\leq l + k\cdot\frac{B}{A} < C$, also gibt es (unabhängig vom Wert von $k$) genau $C$ mögliche Werte von $l$. Insgesamt haben wir also $AC$ Werte von $u$ in $\mathcal P$ mit $C\tau u\in L$.
Weil $u=0$ keine $C\tau$-Teilungsstelle ist, folgt $|\DIV(C\tau)|=AC-1$.}

\en{For all $u\in \DIV(C\tau)$, it holds eq.~(\ref{\en{EN}ctauu}), which implies $AC\cdot u \in L$ and thus $u\in \DIV(AC)$.
This proves that $\DIV(C\tau)\subset \DIV(AC)$.}%
\de{Für alle $u\in \DIV(C\tau)$ gilt Glg.~(\ref{\en{EN}ctauu}). Hieraus folgt $AC\cdot u \in L$ und somit $u\in \DIV(AC)$.
Somit haben wir auch $\DIV(C\tau)\subset \DIV(AC)$ bewiesen.}
\end{proof}




\begin{lem}\label{\en{EN}lemkonst}
\en{In the notations of Def.~\ref{\en{EN}bemkonv} and with the function $f(z)$ from Lemma~\ref{\en{EN}lemfellip} it holds:}%
\de{In der Notation von Def.~\ref{\en{EN}bemkonv} und mit der Funktion $f(z)$ aus Lemma~\ref{\en{EN}lemfellip} gilt:}
$$g(z):=f(z)-\left(1-\frac{A}{C}\right)\zeta(z)+ \frac{A}{C} \cdot\sum_{u\in \DIV(C)}\zeta(z-u) - \sum_{v\in \DIV(C\tau)}\zeta(z-v)$$
\en{is constant in the whole complex plane.}%
\de{ist konstant in der ganzen komplexen Ebene.}
\end{lem}
\begin{proof}
\en{First we observe that $f(z)$ has poles of order $1$ with residuum $-A/C$ in all $C$-division-points, and poles with residuum $\frac{C\tau}{C\tau}=1$ in all $C\tau$-division-points. In $z=0$, $f(z)$ has a pole with residuum $-\frac A C+1$.}%
\de{Zunächst bemerken wir, dass $f(z)$ Pole erster Ordnung mit Residuum $-A/C$ in allen $C$-Teilungsstellen hat, und Pole mit Residuum $\frac{C\tau}{C\tau}=1$ in allen $C\tau$-Teilungsstellen. In $z=0$ hat $f(z)$ einen Pol mit Residuum $-\frac A C+1$.}
\en{Thus, by definition of $g(z)$, $g(z)$ has no poles and is analytic in the whole complex plane.}%
\de{Aus der Definition von $g(z)$ folgt also, dass $g(z)$ keine Pole hat und in der ganzen komplexen Ebene analytisch ist.}
\en{We will now prove that $g$ is also elliptic. For this we observe that there are $C^2-1$ values in the $u$-summation (Lemma~\ref{\en{EN}lemmm}) and $AC-1$ values in the $v$-summation (Lemma~\ref{\en{EN}lemctau}).
This yields:}%
\de{Wir werden jetzt beweisen, dass $g$ auch elliptisch ist. Hierfür nutzen wir, dass es $C^2-1$ Werte in der $u$-Summation (Lemma~\ref{\en{EN}lemmm}) und $AC-1$ Werte in der $v$-Summation (Lemma~\ref{\en{EN}lemctau}) gibt.
Daraus folgt nämlich:}
\begin{align*}
    &~~~~~~~g(z+\omega_k)-g(z)\\
    &= f(z+\omega_k)-f(z)-\left(1-\frac{A}{C}\right)\cdot \eta_k + \frac A C \cdot (C^2-1)\cdot \eta_k - (AC-1)\cdot \eta_k\\
    &= f(z+\omega_k)-f(z)+\eta_k\cdot\left(-1+\frac{A}{C}+ AC-\frac A C - AC + 1\right) = f(z+\omega_k)-f(z)
\end{align*}
\en{Since $f(z+\omega_k)=f(z)$ (cf.~Lemma~\ref{\en{EN}lemfellip}) we get $g(z+\omega_k)-g(z)=0$. Since all elliptic functions without poles are constant (cf.~first Liouville theorem, Prop.~\ref{\en{EN}liouville1}), the Lemma is proven.}%
\de{Aus Lemma~\ref{\en{EN}lemfellip} folgt $f(z+\omega_k)=f(z)$ und somit $g(z+\omega_k)-g(z)=0$.
Da alle elliptischen Funktionen ohne Pole konstant sind (erster Liouville'scher Satz~\ref{\en{EN}liouville1}), ist das Lemma bewiesen.}
\end{proof}




\begin{lem}\label{\en{EN}lemkapa2}
\en{In the notations of Def.~\ref{\en{EN}bemkonv} and with $\kappa$ from Def.~\ref{\en{EN}defkappa} it holds:}%
\de{In der Notation von Def.~\ref{\en{EN}bemkonv} und mit $\kappa$ aus Def.~\ref{\en{EN}defkappa} gilt:}
$$C\tau\kappa = - \sum_{v\in \DIV(C\tau)}\wp(v)$$
\end{lem}
\begin{proof}
\en{From Def.~\ref{\en{EN}defiwp} we get $\zeta'(z+w)=-\wp(z+w)$ and thus around $z=0$ it holds:
$\zeta(z+w)=\zeta(w)-\wp(w)\cdot z + O(z^2)$ (if $w\notin L$).}%
\de{Aus Def.~\ref{\en{EN}defiwp} folgt $\zeta'(z+w)=-\wp(z+w)$, somit gilt um $z=0$ (falls $w\notin L$):
$\zeta(z+w)=\zeta(w)-\wp(w)\cdot z + O(z^2)$.}
\en{Def.~\ref{\en{EN}defizeta} along with Prop.~\ref{\en{EN}sigmaodd} proves that $\zeta(z)$ is odd:
$\zeta(-z)=\frac{\sigma'(-z)}{\sigma(-z)}=\frac{\sigma'(z)}{-\sigma(z)}=-\zeta(z)$.
In Prop.~\ref{\en{EN}laurentwp}, we proved that the Laurent series of $\wp(z)$ at $z=0$ has no constant term.
From $\wp(z)=-\zeta'(z)$ we deduce that around $z=0$ it holds $\zeta(z)=\frac 1 z + O(z^3)$.}%
\de{Aus Def.~\ref{\en{EN}defizeta} folgt mit Satz~\ref{\en{EN}sigmaodd}, dass $\zeta(z)$ ungerade ist:
$\zeta(-z)=\frac{\sigma'(-z)}{\sigma(-z)}=\frac{\sigma'(z)}{-\sigma(z)}=-\zeta(z)$.
In Satz~\ref{\en{EN}laurentwp} haben wir bewiesen, dass der konstante Term der Laurentreihe von $\wp(z)$ bei $z=0$ verschwindet.
Mit $\wp(z)=-\zeta'(z)$ folgt dann, dass um $z=0$ gilt: $\zeta(z)=\frac 1 z + O(z^3)$}

\en{Thus we obtain around $z=0$:}\de{Somit folgt um $z=0$:}
\begin{align*}
    g(z) =& -\frac{A}{Cz} + \frac{C\tau}{C\tau z} + C \tau \kappa z - \left(1-\frac A C\right)\cdot \frac 1 z
    + \frac A C \cdot\sum_{u\in \DIV(C)}\left(\zeta(-u)-\wp(-u)\cdot z\right)\\
    &- \sum_{v\in \DIV(C\tau)} \left( \zeta(-v) - \wp(-v)\cdot z \right) + O(z^2)
\end{align*}
\en{Since $g(z)$ is constant (cf.~Lemma~\ref{\en{EN}lemkonst}), the coefficient of $z$ in this Laurent series expansion is zero:}%
\de{Da $g(z)$ konstant ist (vgl.~Lemma~\ref{\en{EN}lemkonst}), verschwindet der Koeffizient vor $z$ in dieser Laurentreihe:}
\begin{align*}
    0 &= C \tau \kappa - \frac A C \cdot\sum_{u\in \DIV(C)}\wp(u) + \sum_{v\in \DIV(C\tau)} \wp(v)
\end{align*}
\en{Thm.~\ref{\en{EN}theowpsum} tells that $\sum_{u\in \DIV(C)}\wp(u)=0$.
It remains $0 = C\tau\kappa + \sum_{v\in \DIV(C\tau)}\wp(v)$ which proves the Lemma.}%
\de{Thm.~\ref{\en{EN}theowpsum} besagt, dass $\sum_{u\in \DIV(C)}\wp(u)=0$.
Es verbleibt $0 = C\tau\kappa + \sum_{v\in \DIV(C\tau)}\wp(v)$, womit das Lemma bewiesen ist.}
\end{proof}

\begin{thm}\label{\en{EN}lemkapppppa}
\en{In the notations of Def.~\ref{\en{EN}bemkonv}, it holds:}%
\de{In der Notation von Def.~\ref{\en{EN}bemkonv} gilt:}
$$\sqrt{D}\cdot E_2^*(\tau)\cdot\frac{\pi^2}{3\omega_1^2} = \frac{\sum_{v\in \DIV(C\tau)}\wp(v)}{C\tau}$$
\end{thm}
\begin{proof}
\en{We have proved two representations of $\kappa$ in Lemma~\ref{\en{EN}lemkapa1} and Lemma~\ref{\en{EN}lemkapa2}. Equating these yields Prop.~\ref{\en{EN}lemkapppppa}.}%
\de{Wir haben in Lemma~\ref{\en{EN}lemkapa1} und~\ref{\en{EN}lemkapa2} zwei Darstellungen von $\kappa$ bewiesen. Diese gleichzusetzen liefert Satz~\ref{\en{EN}lemkapppppa}.}
\end{proof}


\begin{theo}[\en{Remainder of Prop.}\de{Rest des Satzes}~\ref{\en{EN}satzezweistern}]\label{\en{EN}theoezweisternrest}
\en{Let $\eta(\tau)$ denote the Dedekind $\eta$-function with $1728\eta^{24}=E_4^3-E_6^2$.
Then, in the notations of Def.~\ref{\en{EN}bemkonv},}%
\de{$\eta(\tau)$ bezeichne die Dedekind'sche $\eta$-Funktion mit $1728\eta^{24}=E_4^3-E_6^2$.
Dann gilt mit den Notationen aus Def.~\ref{\en{EN}bemkonv}, dass }%
\begin{align*}
    \sqrt{D}\cdot\frac{E_2^*(\tau)}{\eta^4(\tau)}\cdot(AC)^2
\end{align*}
\en{is an algebraic integer if it holds $C\tau^2+B\tau+A=0$.}%
\de{ganzalgebraisch ist, falls $C\tau^2+B\tau+A=0$ gilt.}
\end{theo}

\begin{proof}\label{\en{EN}proofzweistern}
\en{For $\tau\in CM$ we denote the lattice}%
\de{Für $\tau\in CM$ bezeichnen wir das Gitter}
$\hat L:=\frac{\pi}{\sqrt{3}}\cdot\eta(\tau)^2\cdot L_\tau$.

\en{Then we have $\hat L=a\cdot L_\tau$ with $a=\frac{\pi}{\sqrt{3}}\cdot\eta(\tau)^2$ and thus Prop.~\ref{\en{EN}trafog23} and Thm.~\ref{\en{EN}fouriertheorem} yield:}%
\de{Dann gilt $\hat L=a\cdot L_\tau$ mit $a=\frac{\pi}{\sqrt{3}}\cdot\eta(\tau)^2$.
Somit folgt aus Satz~\ref{\en{EN}trafog23} und Thm.~\ref{\en{EN}fouriertheorem}:}
\begin{align*}
    h_2 := \frac 1 4 g_2(\hat L) &= \frac 1 4 \cdot a^{-4} \cdot g_2(L_\tau)=\frac 1 4 \cdot \left(\frac {\sqrt{3}}\pi\right)^4\cdot\frac{\frac 4 3\cdot\pi^4\cdot E_4(\tau)}{\eta(\tau)^8} = 3\cdot \frac{E_4(\tau)}{\eta(\tau)^8}\\
    h_3 := \frac 1 4 g_3(\hat L) &= \frac 1 4 \cdot a^{-6} \cdot g_3(L_\tau)=\frac 1 4\cdot \left(\frac {\sqrt{3}}\pi\right)^6\cdot\frac{\frac 8 {27}\cdot\pi^6\cdot E_6(\tau)}{\eta(\tau)^{12}} = 2\cdot \frac{E_6(\tau)}{\eta(\tau)^{12}}
\end{align*}
\en{Since we already proved in Prop.~\ref{\en{EN}satzezweistern} on p.~\pageref{\en{EN}satzezweistern} that $\frac{E_4(\tau)}{\eta(\tau)^8}$ and $\frac{E_6(\tau)}{\eta(\tau)^{12}}$ are algebraic integers, both $h_2 := \frac 1 4 g_2(\hat L)$ and $h_3 := \frac 1 4 g_3(\hat L)$ are algebraic integers.}%
\de{In Satz~\ref{\en{EN}satzezweistern} auf S.~\pageref{\en{EN}satzezweistern} haben wir bereits bewiesen, dass $\frac{E_4(\tau)}{\eta(\tau)^8}$ und $\frac{E_6(\tau)}{\eta(\tau)^{12}}$ ganzalgebraisch sind, also sind sowohl $h_2 := \frac 1 4 g_2(\hat L)$ als auch $h_3 := \frac 1 4 g_3(\hat L)$ ganzalgebraisch.}

\en{In Prop.~\ref{\en{EN}lemkapppppa} we proved for any lattice $L$ with complex multiplication:}%
\de{In Satz~\ref{\en{EN}lemkapppppa} haben wir für alle Gitter $L$ mit komplexer Multiplikation bewiesen:}
$$\sqrt{D}\cdot E_2^*(\tau)\cdot\frac{\pi^2}{3\omega_1^2} = \frac{\sum_{v\in \DIV(C\tau)}\wp(v;L)}{C\tau}.$$
\en{For the lattice $\hat L$ we have $\omega_1 = \frac{\pi}{\sqrt{3}}\cdot\eta(\tau)^2$ and thus $\frac{\pi^2}{3\omega_1^2}= \frac{1}{\eta^4(\tau)}$.
This proves that}%
\de{Für das Gitter $\hat L$ gilt $\omega_1 = \frac{\pi}{\sqrt{3}}\cdot\eta(\tau)^2$ und somit $\frac{\pi^2}{3\omega_1^2}= \frac{1}{\eta^4(\tau)}$.
Hieraus folgt}
\begin{align*}
\sqrt{D}\cdot \frac{E_2^*(\tau)}{\eta^4(\tau)}
= \frac{\sum_{v\in \DIV(C\tau)}\wp(v;\hat L)\cdot C\bar\tau}{C\tau\cdot C\bar\tau}
= \frac{\sum_{v\in \DIV(C\tau)}\wp(v;\hat L)\cdot C\bar\tau}{AC}
\end{align*}
\en{where we expanded the fraction with $C\bar\tau$ to obtain $C\tau\cdot C\bar\tau=AC$ in the denominator.}%
\de{wobei wir mit $C\bar\tau$ erweitert haben, um den Nenner von $C\tau\cdot C\bar\tau=AC$ zu erreichen.}\linebreak
\en{If we multiply this by $(AC)^2$, we obtain}%
\de{Durch Multiplikation mit $(AC)^2$ erhalten wir}
\begin{align*}
\sqrt{D}\cdot \frac{E_2^*(\tau)}{\eta^4(\tau)}\cdot (AC)^2 &= \left(\sum_{v\in \DIV(C\tau)}AC\cdot\wp(v;\hat L)\right)\cdot C\bar\tau
\end{align*}
\en{After Lemma~\ref{\en{EN}lemctau} we know that $v\in \DIV(C\tau)$ implies $v\in \DIV(AC)$.
And above we proved that $h_2$ and $h_3$ are algebraic integers, thus by Thm.~\ref{\en{EN}theomwpu} with $m=AC$, the summands $AC\cdot\wp(v;\hat L)$ are algebraic integers for all $v\in \DIV(AC)$ and for all $v\in \DIV(C\tau)$.
Thus $AC\cdot \sum_{v\in \DIV(C\tau)}\wp(v;\hat L)$ is an algebraic integer.\\
Since $C\bar\tau$ is a solution of $x^2+Bx+AC=0$, it is also an algebraic integer.\linebreak
As the product of algebraic integers is an algebraic integer itself, Thm.~\ref{\en{EN}theoezweisternrest} is proven.}%
\de{Wegen Lemma~\ref{\en{EN}lemctau} wissen wir, dass alle $v\in \DIV(C\tau)$ auch in $\DIV(AC)$ liegen.
Außerdem haben wir soeben bewiesen, dass $h_2$ und $h_3$ ganzalgebraisch sind.
Aus Thm.~\ref{\en{EN}theomwpu} folgt nun (mit $m=AC$), dass die Summanden $AC\cdot\wp(v;\hat L)$ für alle $v\in \DIV(AC)$ und somit für alle $v\in \DIV(C\tau)$ ganzalgebraisch sind.

Folglich ist $AC\cdot \sum_{v\in \DIV(C\tau)}\wp(v;\hat L)$ ganzalgebraisch.
Da $C\bar\tau$ eine Lösung von $x^2+Bx+AC=0$ ist, ist es ebenfalls ganzalgebraisch und Thm.~\ref{\en{EN}theoezweisternrest} ist bewiesen.}
\end{proof}


\begin{bem}
\en{The expression $X:=\sqrt{D}\cdot\frac{E_2^*(\tau)}{\eta^4(\tau)}$ is invariant under the transformations $\tau\rightarrow -1/\tau$ and $\tau\rightarrow \tau+N$ (for $N\in \mathbb Z$), but the value of $(AC)^2$ changes.
For example $\tau\rightarrow\tau'=\tau+1$ transforms $\tau^2-\tau+41=0$ into $(\tau'-1)^2-(\tau'-1)+41=0$ or $\tau'^2-3\tau'+43=0$.
Thus Thm.~\ref{\en{EN}theoezweisternrest} tells for this $\tau$, that $X\cdot 41^2$ as well as $X\cdot 43^2$ are algebraic integers.}%
\de{Der Ausdruck $X:=\sqrt{D}\cdot\frac{E_2^*(\tau)}{\eta^4(\tau)}$ ist invariant unter den Transformationen $\tau\rightarrow -1/\tau$ und $\tau\rightarrow \tau+N$ (für $N\in \mathbb Z$). Allerdings ändert sich dabei der Wert von $(AC)^2$.
Auf diese Weise kann man z.B.~mit $\tau\rightarrow\tau'=\tau+1$ aus der Gleichung $\tau^2-\tau+41=0$ die Gleichung $(\tau'-1)^2-(\tau'-1)+41=0$ bzw.~$\tau'^2-3\tau'+43=0$ erzeugen.
Thm.~\ref{\en{EN}theoezweisternrest} besagt also für dieses $\tau$, dass sowohl $X\cdot 41^2$ als auch $X\cdot 43^2$ ganzalgebraisch ist.}

\en{Euclid's algorithm yields $u,v\in\mathbb Z$ with $u\cdot 41^2+v\cdot43^2 = \operatorname{gcd}(41^2;43^2)=1$.
This proves that $X$ itself is an algebraic integer, because}%
\de{Der Euklidische Algorithmus liefert $u,v\in\mathbb Z$ mit $u\cdot 41^2+v\cdot43^2 = \operatorname{ggT}(41^2;43^2)=1$.
Hieraus folgt, dass auch $X$ selbst ganzalgebraisch ist, denn}
$$X=(u\cdot 41^2+v\cdot43^2)\cdot X = u\cdot X\cdot 41^2 + v\cdot X\cdot 43^2$$

\en{By applying this method repeatedly, one can prove (using the factor $\frac 1 2$ from Thm.~\ref{\en{EN}theomwpu}) that for all $\tau\in CM$, the value $2\cdot X$ is an algebraic integer, without the factor $(AC)^2$.

And if $AC$ is odd or if $B$ is even, one can prove with this method that $X$ itself is an algebraic integer (which is stated without proof in \cite[Prop.~5.10.6]{\en{EN}CohenStromberg}).}%
\de{Durch wiederholtes Anwenden dieser Methode kann man (mit dem Faktor $\frac 1 2$ aus Thm.~\ref{\en{EN}theomwpu}) für alle $\tau\in CM$ beweisen, dass $2\cdot X$ ganzalgebraisch ist, ohne den Faktor $(AC)^2$.

Und falls $AC$ ungerade ist oder $B$ gerade ist, kann man hiermit sogar beweisen, dass $X$ ganzalgebraisch ist (das wird in \cite[Prop.~5.10.6]{\en{EN}CohenStromberg} formuliert, dort aber ohne Beweis).}
\end{bem}








